\documentclass[11pt,a4paper]{article}
\usepackage[margin=1in]{geometry}
\usepackage{amsmath,amssymb,amsfonts}
\usepackage{mathtools}
\usepackage{lmodern}
\usepackage[T1]{fontenc}
\usepackage[utf8]{inputenc}
\usepackage{textcomp}
\usepackage{microtype}
\usepackage{parskip}
\usepackage{enumitem}
\usepackage{array}
\usepackage{booktabs}
\usepackage{hyperref}
\hypersetup{hidelinks,
  pdftitle={Paper 031 -- Slope-4 Baryonic Tully-Fisher Relation with Zero Intrinsic Scatter},
  pdfauthor={Allen Proxmire}}
\setlength{\parindent}{0pt}
\setlength{\parskip}{6pt plus 2pt minus 1pt}

\title{\vspace{-1cm}\Large\bfseries Paper 031 --- Slope-4 Baryonic Tully-Fisher Relation with Zero Intrinsic Scatter}
\author{Allen Proxmire}
\date{May 2026}

\begin{document}
\maketitle

\section*{Abstract}

The Event Density (ED) substrate is a 13-primitive generative system. This paper develops the \textbf{baryonic Tully-Fisher relation (BTFR)}:
\[
v_\mathrm{flat}^{\,4} = G \cdot M_b \cdot a_0
\]
as a \textbf{downstream structural consequence} of four prior results in the ED gravity arc: Newton's law with $G = c^3\ell_P^{\,2}/\hbar$ (Paper\_027), the transition acceleration $a_0 = cH_0/(2\pi)$ (Paper\_029), the substrate-level \textbf{ED Combination Rule} $a = \sqrt{a_N \cdot a_0}$ derived in Paper\_030 (conditional on P14, Bilocal Strain Coupling), and the standard circular-orbit kinematic condition $v^2/R = a$. The BTFR slope is \textbf{exactly 4} --- not 3, not 5 --- because the squaring of $v^2$ in $v^4$ is structurally fixed by the circular-orbit condition combined with the multiplicative form of the ECR. The \textbf{intrinsic scatter is zero} in the deep-MOND asymptotic regime ($a_N \ll a_0$) because the radial dependence cancels exactly: $v^2 = R\sqrt{a_N a_0} = \sqrt{GM_b a_0}$ depends only on baryonic mass $M_b$ and the substrate-level transition acceleration $a_0$, not on the radial profile of $M_b(R)$. \textbf{No halo modeling is invoked. No free parameters are introduced.} Empirical BTFR measurements (McGaugh-Lelli-Schombert 2016 and subsequent surveys) are consistent with slope-4 and small residual scatter attributable to observational uncertainty. The paper makes no claim to derive $H_0$ or specific $M_b$ values, no claim that BTFR is forced from nothing, no claim that MOND is correct as a complete theory of galactic dynamics, and no claim about dark-matter particle content. The empirical case for the substrate primitives + P14 --- and for the slope-4 + zero-scatter prediction --- rests on cross-domain reach.

\section{Introduction}

\subsection{What this paper does}

The \textbf{baryonic Tully-Fisher relation (BTFR)} is one of the tightest empirical relations in galactic astrophysics: the asymptotic flat rotation velocity $v_\mathrm{flat}$ of a galaxy is related to its total baryonic mass $M_b$ by a power law with slope very close to 4:
\[
v_\mathrm{flat}^{\,n} \propto M_b, \quad n \approx 4.
\]
Observational data (McGaugh 2012; Lelli, McGaugh, Schombert 2016) give $n = 3.95 \pm 0.08$, consistent with $n = 4$ to within statistical error. The relation holds across galaxies spanning $\sim$4 orders of magnitude in baryonic mass with remarkably small scatter.

In ED's substrate framework, BTFR slope-4 emerges as the \textbf{downstream structural consequence} of:

\begin{enumerate}[noitemsep]
\item Newton's law $a_N = GM/R^2$ (Paper\_027).
\item The transition acceleration $a_0 = cH_0/(2\pi)$ (Paper\_029).
\item The \textbf{ED Combination Rule} $a = \sqrt{a_N \cdot a_0}$ derived in \textbf{Paper\_030} (conditional on \textbf{P14, Bilocal Strain Coupling}).
\item The circular-orbit condition $v^2/R = a$ (standard kinematics).
\end{enumerate}

Combining these yields:
\[
v_\mathrm{flat}^{\,4} = G M_b a_0.
\]
The slope is \emph{exactly 4} --- structurally fixed by the squaring of $v^2$ in the centripetal-acceleration condition combined with the ECR's multiplicative form. The intrinsic scatter is \emph{zero} in the deep-MOND asymptotic regime because the radial profile of $M_b(R)$ cancels exactly. No halo modeling, no free parameters.

\subsection{Why BTFR matters in the ED generative system}

The empirical BTFR is one of the strongest constraints on galactic gravitational physics. Three frameworks attempt to account for it:

\begin{itemize}[noitemsep]
\item \textbf{Dark-matter cosmology:} BTFR slope-4 emerges from tuned halo profile distributions; the tightness is a coincidence of the tuning.
\item \textbf{MOND-as-phenomenology:} BTFR slope-4 is built in by construction.
\item \textbf{ED substrate gravity:} BTFR slope-4 emerges as the unique downstream consequence of the substrate-level cumulative-strain reading (Paper\_027) + cosmic-horizon dipole projection (Paper\_029) + ECR (Paper\_030, conditional on P14). The slope is \emph{derived}, not fit.
\end{itemize}

This matters for three reasons:

\begin{itemize}[noitemsep]
\item \textbf{Predictive content.} ED's substrate framework makes a specific prediction (slope-4, zero intrinsic scatter) without free parameters.
\item \textbf{Cross-domain coherence.} BTFR is the natural product of Papers\_027 + 029 + 030. It completes the gravity arc.
\item \textbf{Falsifiability.} If empirical BTFR slope deviates from 4 (beyond observational uncertainty), or if intrinsic scatter is found in the deep-MOND regime, the substrate-level mechanism is falsified.
\end{itemize}

\subsection{How this fits into the gravity arc}

The ED substrate-gravity arc consists of four papers:

\begin{itemize}[noitemsep]
\item \textbf{Paper\_027:} Newton's $G = c^3\ell_P^{\,2}/\hbar$ from substrate constants.
\item \textbf{Paper\_029:} Transition acceleration $a_0 = cH_0/(2\pi)$ from cosmic-decoupling-surface dipole projection.
\item \textbf{Paper\_030:} \textbf{Full structural derivation of the ED Combination Rule $a = \sqrt{a_N \cdot a_0}$}, given postulate \textbf{P14 (Bilocal Strain Coupling)} and the joint weak-gradient regime.
\item \textbf{Paper\_031 (this paper):} Slope-4 BTFR $v^4 = GM_b a_0$ with zero intrinsic scatter, composing Paper\_027 + Paper\_029 + Paper\_030 + circular-orbit kinematics.
\end{itemize}

\textbf{Repo-structure note.} The ECR derivation lives in \textbf{Paper\_030 as a dedicated paper}, not in this paper. Paper\_031 \emph{uses} the ECR result (Paper\_030 \S6.1) as a structural step in the BTFR derivation; it does \emph{not} re-derive the ECR. The current arc structure separates ECR derivation (Paper\_030) from BTFR application (Paper\_031) for cleaner scope and easier cross-reference.

The four papers together produce the complete ED substrate-level account of galactic gravity: Newton in the strong-acceleration regime, MOND-class transition at $a_0$, multiplicative combination in the joint regime (via P14), and flat rotation curves with slope-4 BTFR in the deep-MOND regime --- all from the postulated substrate primitives + P14.

\section{Primitive Inputs (postulated, not derived in this paper)}

This paper takes the following Event Density (ED) substrate primitives as \textbf{postulated within the ED Generative Primitives System}:

\begin{itemize}[noitemsep]
\item \textbf{P02 (Participation as primitive relation).}
\item \textbf{P03 (Channel + locus indexing; spatial homogeneity).}
\item \textbf{P04 (Bandwidth as non-negative additive scalar).}
\item \textbf{P07 (Channel structure as ontological primitive).}
\item \textbf{P08 (Substrate scale $\ell_\mathrm{ED}$).}
\item \textbf{P11 (Commitment irreversibility).}
\item \textbf{P12 (Stability landscape).} Load-bearing for the substrate-level mechanism producing the ECR's logarithmic cross-term (in Paper\_030).
\item \textbf{P13 (Time homogeneity).}
\item \textbf{P14 (Bilocal Strain Coupling).} Postulate articulated in Paper\_030 \S2. \textbf{Load-bearing in this paper as an inherited dependency:} the ECR result used here (\S3 below) is conditional on P14.
\end{itemize}

\textbf{V1 inheritance (Papers \#18, \#19 / Paper\_093):} finite-width retarded kernel structure inherited as in Papers\_027, 029, 030.

\textbf{Upstream paper dependencies:}

\begin{itemize}[noitemsep]
\item \textbf{Paper\_027 (Newton's $G$):} establishes $a_N = GM/R^2$ with $G = c^3\ell_P^{\,2}/\hbar$.
\item \textbf{Paper\_029 (transition acceleration $a_0$):} establishes $a_0 = cH_0/(2\pi)$.
\item \textbf{Paper\_030 (ECR derivation, conditional on P14):} establishes $a = \sqrt{a_N \cdot a_0}$ in the joint weak-gradient regime. \textbf{Direct upstream input.}
\item \textbf{Standard kinematics:} the circular-orbit condition $v^2/R = a$.
\end{itemize}

\textbf{Empirical inputs:} $c$, $\hbar$, $H_0$, $\ell_P$, per-galaxy $M_b$.

\textbf{No primitive forcing is invoked.}

\section{Newton + $a_0$ + ED Combination Rule}

\subsection{Summary of Paper\_027 and Paper\_029}

From Paper\_027:
\[
a_N = \frac{GM}{R^2}, \qquad G = \frac{c^3 \ell_P^{\,2}}{\hbar}.
\]
From Paper\_029:
\[
a_0 = \frac{cH_0}{2\pi} \approx 1.08 \times 10^{-10}\,\mathrm{m/s^2}.
\]
These are the two upstream acceleration scales.

\subsection{The ECR taken from Paper\_030}

\textbf{The ECR used in this paper is taken from Paper\_030 and is conditional on the bilocal-strain-coupling postulate (P14) introduced there.} This paper does not re-derive the ECR; the derivation lives in Paper\_030 \S\S4--6.

The result of Paper\_030, used here as a structural step:
\[
\boxed{a = \sqrt{a_N \cdot a_0} \qquad \text{(ECR, from Paper\_030, given P14 + joint weak-gradient regime)}}
\]

\textbf{Conditional structure.} The ECR is a downstream structural consequence of:

\begin{itemize}[noitemsep]
\item The substrate primitives P02 + P03 + P04 + P07 + P08 + P11 + P12 + P13.
\item P14 (Bilocal Strain Coupling) --- articulated in Paper\_030 \S2.
\item V1 inheritance (Papers \#18, \#19 / Paper\_093).
\item The joint weak-gradient regime assumption.
\item The upstream results: $G$ from Paper\_027 and $a_0$ from Paper\_029.
\end{itemize}

This paper inherits the ECR as a substrate-derived structural result with the conditional dependency on P14 made explicit. If P14 is structurally challenged (see Paper\_030 \S8.4 honest accounting of P14's postulational status), the ECR --- and consequently the BTFR slope-4 + zero-scatter result of this paper --- is structurally challenged accordingly. The empirical case is therefore joint: BTFR slope-4 tests the substrate framework + Paper\_030's ECR + P14 simultaneously.

\subsection{The three regimes (summary)}

\begin{center}
\begin{tabular}{p{0.20\textwidth}p{0.30\textwidth}p{0.42\textwidth}}
\toprule
\textbf{Regime} & \textbf{Acceleration} & \textbf{Empirical} \\
\midrule
Newtonian ($a_N \gg a_0$) & $a \approx a_N = GM/R^2$ & Solar system, inner galactic regions \\
Joint weak-gradient ($a_N \sim a_0$) & $a \approx \sqrt{a_N \cdot a_0}$ (P14 active) & Galactic outer regions, transition zone \\
Deep-MOND ($a_N \ll a_0$) & $a \approx \sqrt{a_N \cdot a_0}$ (P14 active) & Far outer galactic regions, dwarf galaxies \\
\bottomrule
\end{tabular}
\end{center}

The ECR is valid throughout the joint and deep-MOND regimes (where P14 applies). In the Newtonian regime, $a_N$ dominates and the ECR reduces to $a \approx a_N$ via the regime-restriction structure of P14 (Paper\_030 \S2). The asymptotic flat-rotation-curve behavior occurs in the deep-MOND regime, which is where BTFR slope-4 is most precisely defined and tested.

\section{Circular-Orbit Condition}

\subsection{Newtonian kinematics}

For a test object in a circular orbit at radius $R$:
\[
\frac{v^2}{R} = a
\]
where $v$ is orbital velocity and $a$ is total radial acceleration. Standard kinematics; not substrate-level content.

\subsection{Substituting the ECR}

In the deep-MOND regime where ECR applies (Paper\_030 \S6.1), $a \approx \sqrt{a_N \cdot a_0}$. Substituting:
\[
\frac{v^2}{R} = \sqrt{a_N \cdot a_0} = \sqrt{\frac{GM_b}{R^2} \cdot a_0} = \frac{\sqrt{GM_b a_0}}{R}.
\]
Multiplying both sides by $R$:
\[
v^2 = \sqrt{G M_b a_0}.
\]
Squaring:
\[
\boxed{v^4 = G \cdot M_b \cdot a_0.}
\]
This is the slope-4 baryonic Tully-Fisher relation.

\subsection{What ``flat rotation curve'' means in this framework}

A ``flat rotation curve'' is the empirical observation that, in galaxies at large radii:
\[
v(R) \approx v_\mathrm{flat} \quad \text{for } R \gtrsim R_\mathrm{transition}.
\]
Given the ECR (from Paper\_030, conditional on P14), the circular-orbit kinematics produces constant rotation velocity at large radii, with $v_\mathrm{flat}$ uniquely determined by $M_b$ and $a_0$.

\section{Why the Slope is Exactly 4}

\subsection{Structural derivation of the exponent}

From \S4.2:
\[
v^4 = G M_b a_0 \implies v \propto M_b^{1/4} \implies v^n \propto M_b \text{ requires } n = 4.
\]
The slope is \emph{exactly} 4, structurally, because:

\begin{itemize}[noitemsep]
\item The circular-orbit condition is \emph{quadratic} in velocity: $v^2 = aR$.
\item The ECR introduces $\sqrt{a_N \cdot a_0}$ --- a \emph{one-half-power} of $a_N$.
\item Newton's law gives $a_N \propto M_b/R^2$ --- a \emph{first-power} of $M_b$.
\item Composing: $v^2 = R \cdot \sqrt{(GM_b/R^2) \cdot a_0} = \sqrt{GM_b a_0}$, so $v^4 \propto M_b^{\,1}$.
\end{itemize}

The slope-4 emerges from the joint action of centripetal kinematics + ECR square-root + Newton's law's radial cancelation.

\subsection{Why no other exponent is structurally allowed}

Alternative combination rules give different slopes:

\begin{itemize}[noitemsep]
\item \textbf{Additive} ($a = a_N + a_0$): in deep-MOND, $v^2 = a_0 R$; no $M_b$ dependence; slope undefined.
\item \textbf{Maximum} ($a = \max$): same.
\item \textbf{Quadratic} ($a^2 = a_N^{\,2} + a_0^{\,2}$): in deep-MOND, $a \approx a_0$; same.
\item \textbf{Cube-root} ($a^3 = a_N a_0^{\,2}$): slope-6.
\item \textbf{Multiplicative geometric mean} (ECR, $a = \sqrt{a_N a_0}$): slope-4. $\checkmark$
\end{itemize}

Only the ECR --- which Paper\_030 derives from P14 --- produces slope-4. The empirical slope-4 BTFR therefore tests P14 + Paper\_030's ECR derivation jointly.

\subsection{Slope-4 is dimensional / structural, not empirical-fit}

Slope-4 is \emph{not} a parameter ED tunes. It is a structural consequence of the dimensional structure of Newton's law, the ECR square-root, and circular kinematics --- composed without free parameters.

\section{Zero Intrinsic Scatter}

\subsection{The radial cancelation}

From \S4.2:
\[
v^2 = R \cdot \sqrt{a_N \cdot a_0} = R \cdot \sqrt{\frac{GM_b}{R^2} \cdot a_0} = \sqrt{GM_b a_0}.
\]
\textbf{The radius $R$ cancels exactly.} $v_\mathrm{flat}$ depends only on total baryonic mass $M_b$ and the substrate-level transition acceleration $a_0$. It does not depend on the radial distribution of $M_b(R)$, the radius at which the rotation curve is measured (in the asymptotic deep-MOND regime), or detailed potential structure at intermediate radii.

\subsection{Why this produces zero intrinsic scatter}

Empirical BTFR scatter has two sources:

\begin{itemize}[noitemsep]
\item \textbf{Intrinsic scatter:} physical variation in $v_\mathrm{flat}$-$M_b$ beyond measurement.
\item \textbf{Observational scatter:} measurement uncertainty.
\end{itemize}

In ED's substrate framework, the radial cancelation implies asymptotic $v_\mathrm{flat}$ is uniquely determined by $M_b$, independently of any galaxy-specific structural property. \textbf{Intrinsic scatter is zero}; all observed scatter must be observational.

This is sharp: as observational uncertainties shrink, residual scatter should approach zero in the deep-MOND asymptotic regime.

\subsection{Current empirical state}

McGaugh-Lelli-Schombert SPARC (2016): 175 galaxies, slope $3.95 \pm 0.08$, scatter $\sim$0.1 dex. The question is whether the 0.1-dex scatter is observational (consistent with ED) or intrinsic (inconsistent). Current analyses suggest largely observational; tighter measurements could resolve.

\subsection{The deep-MOND asymptotic qualifier}

ED predicts zero intrinsic scatter \textbf{in the deep-MOND asymptotic regime} ($a_N \ll a_0$). In the joint weak-gradient regime ($a_N \sim a_0$), the ECR transitions between Newtonian and deep-MOND behavior, and some residual scatter is expected --- domain-of-applicability statement, not falsification.

\section{Cross-Domain Context}

\subsection{The gravity arc, completed}

This paper completes the ED substrate-gravity arc:

\begin{itemize}[noitemsep]
\item \textbf{Paper\_027:} Newton's $G$.
\item \textbf{Paper\_029:} $a_0 = cH_0/(2\pi)$.
\item \textbf{Paper\_030:} ECR derivation given P14.
\item \textbf{Paper\_031 (this paper):} Slope-4 BTFR with zero intrinsic scatter.
\end{itemize}

All four papers together produce the substrate-level account of galactic gravity, with no halo modeling and no free parameters beyond $c$, $\hbar$, $H_0$, $\ell_P$, per-galaxy $M_b$.

\subsection{Interaction with V5 and cross-domain invariants}

The substrate-gravity arc uses V1 (single-chain vacuum kernel) as the dominant kernel. V5 (Paper \#20 / Paper\_090) operates in different contexts (soft-matter, BH, entanglement). The kernel hierarchy supplies the structural machinery for cross-locus interaction phenomena across ED.

\subsection{Cross-checks across the gravity arc}

BTFR slope-4 depends on:

\begin{itemize}[noitemsep]
\item Paper\_027's $G$ identification.
\item Paper\_029's $a_0$ identification.
\item Paper\_030's ECR derivation + P14.
\item Standard circular-orbit kinematics.
\end{itemize}

Empirical BTFR data --- slope-4 + small scatter --- simultaneously tests all of these. Cross-domain coherence: BTFR is the load-bearing testable consequence of the gravity arc.

\section{What This Paper Does NOT Claim}

\subsection{No derivation of $H_0$}

$H_0$ enters via $a_0 = cH_0/(2\pi)$ (Paper\_029). Empirical cosmological content.

\subsection{No derivation of $M_b$ values}

$M_b$ is astrophysical measurement, value-layer empirical.

\subsection{No claim that BTFR is forced from nothing}

Conditional on postulated primitives + V1 inheritance + P14 (via Paper\_030) + Paper\_027 + Paper\_029 results.

\subsection{No claim that MOND is correct as a complete theory}

ED supplies a substrate mechanism for the empirical phenomenology MOND describes.

\subsection{No claim about dark-matter distributions}

ED's framework operates without halo modeling. Does not refute dark matter as particle content.

\subsection{No claim of MOND-interpolation-function derivation}

ED produces the deep-MOND asymptotic relation + ECR transition (Paper\_030); not the full Milgrom $\mu(x)$.

\subsection{No claim of universal validity across all galaxy classes}

Prediction restricted to the deep-MOND asymptotic regime.

\subsection{No claim of cosmological-evolution prediction}

If $H_0$ varies with cosmic time, $a_0$ varies; BTFR coefficient varies. ED does not derive $H_0$ evolution.

\subsection{No claim that the ECR derivation is contained in this paper}

The ECR derivation is in \textbf{Paper\_030}, not this paper. This paper uses the ECR result. The ECR's structural status (conditional on P14 --- see Paper\_030 \S8.4) flows into this paper as inherited.

\section{Falsification Criteria}

\subsection{Deep-MOND intrinsic scatter is non-zero}

If tighter measurements reduce observational uncertainty but leave residual intrinsic scatter, the radial-cancelation prediction is falsified.

\subsection{Slope deviates from 4}

If empirical slope is significantly distinct from 4 beyond observational uncertainty, the substrate-level mechanism is falsified.

\subsection{$a_0$ does not match $cH_0/(2\pi)$}

If empirical $a_0 \neq cH_0/(2\pi)$ at resolved precision, Paper\_029's mechanism is falsified, BTFR coefficient fails.

\subsection{Multiplicative-combination-rule failure}

If transition-regime rotation curves require non-multiplicative combination, the ECR (Paper\_030, P14) is falsified, BTFR slope-4 with it.

\subsection{Halo modeling required to reproduce BTFR}

If empirical BTFR cannot be reproduced without halo modeling at high precision, the substrate framework's halo-free claim is falsified.

\subsection{Outer-galactic rotation curves do not follow ECR}

Systematic deviations from $v(R) \approx (GM_b(<R) a_0)^{1/4}$ falsify the substrate framework's BTFR prediction.

\subsection{P14 falsified independently}

If P14 (Paper\_030's bilocal-strain-coupling postulate) is independently falsified (e.g., via cross-domain inconsistency or laboratory substrate-probing), Paper\_030's ECR derivation fails, BTFR slope-4 follows.

\section{Position Statement}

The baryonic Tully-Fisher relation with slope \textbf{exactly 4} and \textbf{zero intrinsic scatter} in the deep-MOND asymptotic regime is \textbf{not} a phenomenological coincidence, an empirical fit, or a tuned halo result. In the ED Generative Primitives System, BTFR is a downstream structural identification:
\[
v_\mathrm{flat}^{\,4} = G \cdot M_b \cdot a_0
\]
where:

\begin{itemize}[noitemsep]
\item $G = c^3 \ell_P^{\,2} / \hbar$ (Paper\_027).
\item $a_0 = cH_0/(2\pi)$ (Paper\_029).
\item The slope-4 form is structurally fixed by circular-orbit kinematics + ECR (Paper\_030, conditional on P14) + Newton's law's radial cancelation.
\item Zero intrinsic scatter follows from radial cancelation: $v^2 = R\sqrt{a_N a_0} = \sqrt{GM_b a_0}$ depends only on $M_b$.
\end{itemize}

\textbf{No halo modeling. No free parameters.} Empirically: slope $\approx 3.95 \pm 0.08$, scatter $\sim$0.1 dex.

What this paper claims:

\begin{itemize}[noitemsep]
\item Given the postulated primitives + V1 inheritance + Paper\_027 + Paper\_029 + \textbf{Paper\_030 (ECR, conditional on P14)} + standard kinematics, slope-4 BTFR with zero intrinsic scatter is the unique downstream structural prediction.
\item The substrate-level multiplicative geometric mean (ECR from Paper\_030) is the unique combination rule producing slope-4.
\item Radial cancelation in the deep-MOND regime is the structural mechanism for zero intrinsic scatter.
\item The substrate framework supplies a mechanism for the empirically observed BTFR that neither MOND-as-phenomenology nor dark-matter-as-particle provides.
\end{itemize}

What this paper does \emph{not} claim:

\begin{itemize}[noitemsep]
\item $H_0$ and cosmology not derived.
\item $M_b$ values are value-layer.
\item BTFR not forced from nothing; conditional on the substrate ontology + P14.
\item MOND not endorsed as complete theory.
\item Dark matter not refuted as particle content.
\item Full MOND interpolation function not derived.
\item The ECR derivation is not contained in this paper; it lives in Paper\_030.
\end{itemize}

\textbf{Series context.} Paper\_031 of the gravity arc. The arc is complete with this paper:

\begin{itemize}[noitemsep]
\item \textbf{Paper\_027:} Newton's $G$.
\item \textbf{Paper\_029:} $a_0$.
\item \textbf{Paper\_030:} ECR (given P14).
\item \textbf{Paper\_031 (this paper):} BTFR slope-4 + zero scatter.
\end{itemize}

\section*{Appendix: Cross-References and Glossary}

\subsection*{Cross-references}

\begin{itemize}[noitemsep]
\item \textbf{Position paper:} \texttt{paper\_ED\_Framework\_13\_Primitive\_Generative\_System.md}.
\item \textbf{Primitive load-bearing audit:} \texttt{PRIMITIVE\_LOAD\_BEARING\_AUDIT.md} --- P12 audit; P14 audit pending in Paper\_030.
\item \textbf{Paper\_027 (Newton's $G$):} establishes $G = c^3\ell_P^{\,2}/\hbar$.
\item \textbf{Paper\_029 (transition acceleration $a_0$):} establishes $a_0 = cH_0/(2\pi)$.
\item \textbf{Paper\_030 (ECR, conditional on P14):} establishes $a = \sqrt{a_N \cdot a_0}$. \textbf{Direct upstream input.}
\item \textbf{Paper \#18, \#19 / Paper\_093 (V1 kernel):} finite-width retarded inheritance.
\item \textbf{Paper \#20 / Paper\_090 (V5):} cross-domain context \S7.2.
\end{itemize}

\subsection*{Glossary}

\begin{itemize}[noitemsep]
\item \textbf{Baryonic Tully-Fisher Relation (BTFR).} $v_\mathrm{flat}^{\,4} = GM_b a_0$. Asymptotic flat-rotation-velocity-baryonic-mass relation.
\item \textbf{Slope-4.} Exponent $n = 4$ in $v^n \propto M_b$. Empirically $n \approx 3.95 \pm 0.08$.
\item \textbf{Zero intrinsic scatter.} Substrate-level prediction that all empirical BTFR scatter is observational.
\item \textbf{ED Combination Rule (ECR).} $a = \sqrt{a_N \cdot a_0}$. Derived in \textbf{Paper\_030} conditional on P14 + joint weak-gradient regime. Used here as a structural step.
\item \textbf{P14 (Bilocal Strain Coupling).} Postulate articulated in Paper\_030 \S2. Load-bearing upstream dependency of this paper's BTFR result.
\item \textbf{Joint weak-gradient regime.} $a_N \sim a_0$. Where P14 applies and ECR holds.
\item \textbf{Deep-MOND regime.} $a_N \ll a_0$. ECR dominates; flat rotation curves; BTFR slope-4 applies.
\item \textbf{Circular-orbit condition.} $v^2/R = a$. Standard kinematics.
\item \textbf{Radial cancelation.} Structural mechanism by which $R$-dependence cancels in $v^2 = R\sqrt{a_N a_0} = \sqrt{GM_b a_0}$, giving zero intrinsic scatter.
\item \textbf{Asymptotic flat rotation curve.} $v(R) \approx v_\mathrm{flat}$ for $R \gtrsim R_\mathrm{transition}$.
\end{itemize}

\end{document}
